\section{03 Linux GPIO}

\subsection{GPIO Registers (BCM2711)}

%% >>> IMAGE (slide 3): RPi4 40-pin header pinout <<<
%% >>> IMAGE (slide 6): BCM2711 block diagram (ARM physical 0x0FE200000 vs legacy 0x7E200000, AXI bridge) <<<

\begin{concept}{GPIO \& its Registers}
    The BCM2711 provides \important{58 GPIOs}. Each pin can be GPIO or an
    \important{alternate function} (UART, I2C, SPI, \dots). Controlled via register banks:
    \begin{itemize}
        \item \important{GPFSEL} -- function select \emph{and} direction (in/out/alt)
        \item \important{GPSET} -- set output high, \important{GPCLR} -- set output low
        \item \important{GPLEV} -- read input level (high/low)
    \end{itemize}
\end{concept}

\begin{definition}{GPFSEL -- Function Select}
    3 bits per pin, \important{10 GPIOs per 32-bit register} (GPFSEL0\dots5).
    \begin{itemize}
        \item \texttt{000} input, \texttt{001} output
        \item \texttt{100/101/110/111} alt fn 0--3, \texttt{011/010} alt fn 4--5
    \end{itemize}
    LSB of a pin's field:
\begin{lstlisting}[language=C, style=basesmol, aboveskip=2pt, belowskip=2pt]
#define FSELX_OFFSET(pin) (pin % 10 * 3)  // 3 bits per GPIO
\end{lstlisting}
\end{definition}

\mult{2}

\begin{definition}{SET / CLR / LEV}
    Each is a bitfield, 1 bit per GPIO (banks 0 \& 1 cover 0--57):
    \begin{itemize}
        \item write 1 to \texttt{GPSETn} bit $\to$ pin high
        \item write 1 to \texttt{GPCLRn} bit $\to$ pin low
        \item read \texttt{GPLEVn} bit $\to$ current input level
    \end{itemize}
\end{definition}

\begin{formula}{Register Offsets (base 0x7e200000)}
    \begin{tabular}{l|l}
        \texttt{0x00..0x14} & GPFSEL0\dots5 \\
        \texttt{0x1c / 0x20} & GPSET0 / GPSET1 \\
        \texttt{0x28 / 0x2c} & GPCLR0 / GPCLR1 \\
        \texttt{0x34 / 0x38} & GPLEV0 / GPLEV1 \\
    \end{tabular}
    All accesses are 32-bit.
\end{formula}

\multend

\subsection{GPIO Access Methods}

\begin{concept}{Two Ways to Reach the GPIO Registers}
    \begin{itemize}
        \item \important{Method 1 -- Pseudo filesystem} (\texttt{/sys/class/gpio}):
              needs a GPIO driver, \important{portable}, less hassle with addresses,
              but \emph{slow}
        \item \important{Method 2 -- Memory mapped} (\texttt{/dev/mem} + \texttt{mmap}):
              direct pointer access to peripheral addresses, need the exact address,
              but \important{much faster}
    \end{itemize}
\end{concept}

\begin{example2}{Exam Matching: GPIO access methods (slide 14)}
    Match each description (left) to the right term:
    \begin{enumerate}
        \item Page-Offset
        \item Translates a physical address into a virtual one
        \item Where the physical address of the GPIO register is declared
        \item LED toggled via file manipulation under \texttt{/dev/mem}
        \item LED toggled via file manipulation under \texttt{/sys/class/gpio}
    \end{enumerate}
    Terms: \emph{size of the virtual MMU page} $\cdot$ \texttt{mmap(...)} $\cdot$
    Device Tree $\cdot$ Memory-Map manipulation $\cdot$ Pseudo-FS manipulation
    \tcblower
    \important{Matches:}
    \begin{enumerate}
        \item $\to$ size of the virtual MMU page
        \item $\to$ \texttt{mmap(...)}
        \item $\to$ Device Tree
        \item $\to$ GPIO manipulation via \important{Memory Map}
        \item $\to$ GPIO manipulation via \important{Pseudo File System}
    \end{enumerate}
\end{example2}

\subsection{Method 1: Pseudo Filesystem (sysfs)}

\begin{definition}{sysfs GPIO File Structure}
    Kernel config: \emph{Device drivers $\to$ GPIO Support $\to$ sysfs interface}.
    \begin{itemize}
        \item \texttt{/sys/class/gpio/export}, \texttt{.../unexport} -- (de)activate a GPIO
        \item \texttt{.../gpioXX/direction} -- \texttt{"in"} or \texttt{"out"}
        \item \texttt{.../gpioXX/value} -- \texttt{"1"} or \texttt{"0"}
        \item \texttt{.../gpioXX/edge} -- \texttt{none}/\texttt{rising}/\texttt{falling}/\texttt{both}
    \end{itemize}
\end{definition}

\begin{corollary}{Mind the GPIO Offset! (slide 18--19)}
    On the InES RPi-Hat: LED4=GPIO12, LED5=GPIO13, Key0=GPIO22, Key1=GPIO23.
    But sysfs adds a \important{base offset of 512} (\texttt{gpiochip512},
    label \texttt{pinctrl-bcm2711}). So in sysfs:
    \important{sysfs\_number = 512 + pin} (e.g. GPIO12 $\to$ 524, GPIO22 $\to$ 534).
\end{corollary}

\begin{KR}{Control a GPIO via sysfs}
    \textbf{LED (output), pin 12 $\to$ 524:}
\begin{lstlisting}[language=bash, style=basesmol, aboveskip=2pt, belowskip=2pt]
echo 524 > /sys/class/gpio/export
echo out > /sys/class/gpio/gpio524/direction
echo 1   > /sys/class/gpio/gpio524/value   # on
echo 0   > /sys/class/gpio/gpio524/value   # off
echo 524 > /sys/class/gpio/unexport        # free
\end{lstlisting}
    \textbf{Button (input), pin 22 $\to$ 534:}
\begin{lstlisting}[language=bash, style=basesmol, aboveskip=2pt, belowskip=2pt]
echo 534 > /sys/class/gpio/export
echo in  > /sys/class/gpio/gpio534/direction
cat /sys/class/gpio/gpio534/value          # read
echo 534 > /sys/class/gpio/unexport
\end{lstlisting}
    \paragraph{conclusion}
    Always add the \important{+512 offset} when converting a board GPIO number.
\end{KR}

\begin{examplecode}{sysfs GPIO from C (slides 23--25)}
\begin{lstlisting}[language=C, style=basesmol, aboveskip=2pt, belowskip=2pt]
// export a gpio
FILE *fp = fopen("/sys/class/gpio/export", "w");
fprintf(fp, "%i", gpio); fflush(fp); fclose(fp);

// set value
void fs_gpio_set(int gpio, int val) {
char str[100];
sprintf(str, "/sys/class/gpio/gpio%i/value", gpio);
FILE *fp = fopen(str, "w");
fprintf(fp, "%i", val); fflush(fp); fclose(fp);
}
// read value
int fs_gpio_get(int gpio) {
char str[100]; int ret;
sprintf(str, "/sys/class/gpio/gpio%i/value", gpio);
FILE *fp = fopen(str, "r");
fread(str, 101, 1, fp);
sscanf(str, "%i", &ret); fclose(fp);
return ret;
}
\end{lstlisting}
\end{examplecode}

\subsection{Method 2: Memory-Mapped GPIO}

%% >>> IMAGE (slide 29 / 33): virtual vs physical address space, mmap() maps GPIO_BASE_ADDR (page-aligned) to virtual_gpio_base; legacy bridge 35-bit <-> 32-bit <<<

\begin{concept}{MMU \& mmap}
    The CPU sees \important{virtual} addresses; hardware lives at \important{physical}
    addresses. \texttt{/dev/mem} exposes physical memory; \texttt{mmap()} returns a
    \important{virtual} pointer to the (page-aligned) GPIO base. The program then
    works purely with virtual addresses; the MMU translates them.
\begin{lstlisting}[language=C, style=basesmol, aboveskip=2pt, belowskip=2pt]
void *mmap(addr, len, prot, flags, fd, base_addr);
\end{lstlisting}
\end{concept}

\begin{corollary}{Legacy Bridging}
    The BCM2711 uses \important{35} address lines, but older peripherals still use
    the legacy \important{32}-bit map. A hardware \important{bridge} adapts new$\to$old
    addresses; a \texttt{ranges} section in the \emph{device tree} tells Linux about
    the mapping (e.g. ARM \texttt{0xFE200000} $\leftrightarrow$ legacy \texttt{0x7E200000}).
\end{corollary}

\begin{KR}{Memory-Mapped GPIO Access}
    \textbf{Step 1 -- map the base} (open \texttt{/dev/mem}, \texttt{mmap} one page):
\begin{lstlisting}[language=C, style=basesmol, aboveskip=2pt, belowskip=2pt]
void *virtual_gpio_base;             // global pointer
int mmap_virtual_base() {
int m_mfd;
if ((m_mfd = open("/dev/mem", O_RDWR)) < 0)
return m_mfd;                      // FAIL
virtual_gpio_base = mmap(NULL, sysconf(_SC_PAGE_SIZE),
PROT_READ|PROT_WRITE, MAP_SHARED, m_mfd, GPIO_BASE_ADDR);
close(m_mfd);
if (virtual_gpio_base == MAP_FAILED) return errno;
return 0;
}
\end{lstlisting}
    \textbf{Step 2 -- set/clear} a pin (write the bit into SET0/CLR0):
\begin{lstlisting}[language=C, style=basesmol, aboveskip=2pt, belowskip=2pt]
#define GPIO_SET0 0x1c
#define GPIO_CLR0 0x28
void mmap_gpio_set(int gpio, int value) {
uint32_t *reg = value
? (uint32_t*)(virtual_gpio_base + GPIO_SET0)
: (uint32_t*)(virtual_gpio_base + GPIO_CLR0);
*reg = (0x1 << gpio);
}
\end{lstlisting}
    \textbf{Step 3 -- set direction} (read-modify-write the FSEL field):
\begin{lstlisting}[language=C, style=basesmol, aboveskip=2pt, belowskip=2pt]
#define GPIO_FSEL1 0x04
#define GPIO_FSEL2 0x08
void mmap_gpio_direction(int gpio, int dir) {
uint32_t *reg;
if (gpio <= 9)        reg = (uint32_t*) virtual_gpio_base;
else if (gpio < 20)   reg = (uint32_t*)(virtual_gpio_base + GPIO_FSEL1);
else                  reg = (uint32_t*)(virtual_gpio_base + GPIO_FSEL2);
*reg &= ~(0x7 << FSELX_OFFSET(gpio));   // clear field
*reg |=  (dir << FSELX_OFFSET(gpio));   // set direction
}
\end{lstlisting}
\end{KR}

\subsection{Device Tree}

\begin{concept}{Device Tree -- Purpose}
    A data structure describing the \important{hardware} to Linux (especially
    \important{non-discoverable} hardware). Tells the kernel:
    \begin{itemize}
        \item physical addresses \& sizes of memory/peripheral registers
        \item where interrupts and GPIO/peripheral pins are connected
        \item which alternate functions and \important{which driver} to load (+ enable it)
    \end{itemize}
\end{concept}

\mult{2}

\begin{definition}{Principle \& Files}
    Non-discoverable HW is either coded in C in the kernel \emph{or} described in a
    device tree:
    \begin{itemize}
        \item \texttt{.dts} -- device tree \important{source}
        \item \texttt{.dtb} -- compiled \important{blob} used by the kernel
        \item \texttt{dtc} -- device tree compiler
    \end{itemize}
    The bootloader loads \important{both} kernel image and \texttt{.dtb} before starting.
\end{definition}

\begin{definition}{Mechanism}
    \begin{enumerate}
        \item Device tree is \important{parsed}
        \item a \important{compatible} driver in the kernel is searched \& started
        \item the driver gets additional \important{properties} from the device tree
    \end{enumerate}
\end{definition}

\multend

\begin{definition}{Structure}
    Tree of \texttt{node@unit-address \{ property = value; \dots \}}; nodes can hold
    child nodes; a \important{phandle} (\texttt{<\&label>}) references another node.
    Values: strings, \texttt{<cells>} (32-bit), or \texttt{[bytes]}.
\end{definition}

\begin{example2}{Exam Quiz: Device Tree (slide 42)}
    \textbf{Q1 -- What is a device tree?}
    \begin{enumerate}
        \item Definition of the processor architecture
        \item Data structure about hardware properties
        \item File structure of the filesystem (tree)
        \item Holds info about hardware that can't be auto-detected
    \end{enumerate}
    \textbf{Q2 -- What applies to the term device tree?}
    \begin{enumerate}
        \item The same compiled kernel can support different HW configurations
        \item Holds the virtual addresses of the machine hardware
        \item Holds the device drivers of the machine hardware
        \item Holds a list \& addresses of "non-discoverable" hardware
    \end{enumerate}
    \tcblower
    \important{Q1: (2) and (4).} (1) wrong scope, (3) confuses it with the filesystem.

    \important{Q2: (1) and (4).} (2) false -- it holds \emph{physical} addresses;
    (3) false -- it \emph{references} which driver, it doesn't contain drivers.
\end{example2}

\begin{examplecode}{GPIO node in bcm2711-rpi-4b.dts (slide 47)}
\begin{lstlisting}[language=C, style=basesmol, aboveskip=2pt, belowskip=2pt]
gpio0: gpio@0x7e200000 {
compatible = "brcm,bcm2711-gpio";  // -> find driver
reg = <0x7e200000 0xb4>;           // base addr, size
interrupts = <0x00 0x71 0x04 ...>;
gpio-controller; #gpio-cells = <2>;
interrupt-controller;
pinctrl-names = "default";
phandle = <0x0f>;
i2c1_gpio2 { brcm,pins = <0x02 0x03>;  // GPIO 2 & 3
brcm,function = <0x04>; };            // alt fn 0 = I2C
};
\end{lstlisting}
\end{examplecode}

\begin{example2}{Exam Exercise: Read a Device Tree Node (slide 49)}
\begin{lstlisting}[language=C, style=basesmol, aboveskip=2pt, belowskip=2pt]
i2c@7e804000 {
compatible = "brcm,bcm2711-i2c\0brcm,bcm2835-i2c";
reg = <0x7e804000 0x1000>;
interrupts = <0x00 0x75 0x04>;
status = "okay";
clock-frequency = <0x186a0>;
};
\end{lstlisting}
    \tcblower
    \begin{itemize}
        \item Device? \important{I2C}
        \item Address? \texttt{0x7e804000} -- \important{legacy} view (\texttt{0x7e\dots})
        \item Size? \texttt{0x1000} bytes
        \item Compatible driver? \texttt{brcm,bcm2711-i2c} (the \texttt{\textbackslash0}
              separates a 2nd fallback driver \texttt{bcm2835-i2c})
        \item Driver enabled? \important{yes} (\texttt{status = "okay"}; \texttt{"disabled"}
              $\to$ not loaded at boot)
    \end{itemize}
\end{example2}

\begin{KR}{Reading a Device Tree Node}
    \begin{itemize}
        \item node name before \texttt{@} = \important{device}; after \texttt{@} = address
        \item \texttt{reg = <addr size>} $\to$ register base + length
        \item address \texttt{0x7e\dots} = legacy view, \texttt{0xfe\dots} = ARM view
        \item \texttt{compatible} = driver name(s); \texttt{\textbackslash0} separates fallbacks
        \item \texttt{status = "okay"} loaded, \texttt{"disabled"} not loaded
    \end{itemize}
\end{KR}

\subsection{I2C Bus}

%% >>> IMAGE (slide 51): I2C read/write cycle timing (S, slave addr, R/W, A, DATA, P) <<<

\mult{2}

\begin{concept}{I2C Protocol}
    \begin{itemize}
        \item \important{7-bit} slave address + $R/\overline{W}$ bit, then \important{8-bit}
              data bytes
        \item each byte \important{ACK}'d (SDA low = ACK, high = NACK)
        \item \important{S} = start condition, \important{P} = stop condition
    \end{itemize}
\end{concept}

\begin{definition}{MCP23008 I2C Expander}
    Key registers:
    \begin{itemize}
        \item \texttt{IODIR (0x00)} -- 0 = output, 1 = input
        \item \texttt{GPIO (0x09)} -- read pin levels / write outputs
        \item \texttt{OLAT (0x0A)} -- output latch
        \item \texttt{GPPU (0x06)} -- pull-ups
    \end{itemize}
    sysfs maps its GPIOs as \important{GPIO578\dots585}.
\end{definition}

\multend

\begin{KR}{Inspect \& Drive I2C with i2c-tools}
\begin{lstlisting}[language=bash, style=basesmol, aboveskip=2pt, belowskip=2pt]
lsmod                       # list active drivers
rmmod pinctrl_mcp23s08_i2c  # free bus if a driver holds it
i2cdetect -l                # list buses (e.g. i2c-1)
i2cdetect -y 1              # scan bus 1 (UU=busy, 20=device@0x20)
i2cdump  -y 1 0x20          # dump all regs of device 0x20
i2cset   -y 1 0x20 0x0 0xf0 # write reg 0x00 = 0xf0 (IODIR)
i2cset   -y 1 0x20 0xa 0x0f # write reg 0x0a = 0x0f (OLAT)
\end{lstlisting}
    Syntax: \texttt{i2cset [-y] bus chip-addr data-addr [value]}.
\end{KR}

\subsection{I2C-dev (C API)}

%% >>> IMAGE (slide 69): write vs read sequence to MCP23008 (write reg addr, then read data) <<<

\begin{KR}{Access an I2C Device via i2c-dev}
    \textbf{Step 1 -- open the bus} and select the slave:
\begin{lstlisting}[language=C, style=basesmol, aboveskip=2pt, belowskip=2pt]
#include <linux/i2c-dev.h>
#include <sys/ioctl.h>
int fd = open("/dev/i2c-1", O_RDWR);     // bus file
if (fd < 0) { perror("open i2c"); exit(1); }
int addr = 0x20;
if (ioctl(fd, I2C_SLAVE, addr) < 0)      // pick slave addr
{ printf("no slave\n"); exit(1); }
\end{lstlisting}
    \textbf{Step 2 -- write} (register address + data):
\begin{lstlisting}[language=C, style=basesmol, aboveskip=2pt, belowskip=2pt]
uint8_t wr_buf[2] = { reg_addr, reg_data };
if (write(fd, wr_buf, 2) != 2) printf("write failed\n");
\end{lstlisting}
    \textbf{Step 3 -- read} (write the reg address, then read the data byte):
\begin{lstlisting}[language=C, style=basesmol, aboveskip=2pt, belowskip=2pt]
uint8_t wr_buf[1] = { reg_addr }, rd_buf[1];
write(fd, wr_buf, 1);                    // point to register
if (read(fd, rd_buf, 1) == 1) data = rd_buf[0];
\end{lstlisting}
    \important{Note:} the bus cannot be opened if another driver already occupies it.
\end{KR}

\begin{remark}
    \important{Lab 3--4 goal:} test file-based GPIO and include it in the CLI,
    program memory-mapped GPIO with \texttt{mmap()}, and access the MCP23008 I2C
    expander both with i2c-tools and with i2c-dev.
\end{remark}

\subsection{Sample Exam (2019)}

\begin{example2}{Exam: GPIO via Device Tree \& sysfs (Task 3)}
    \begin{itemize}
        \item driver name? $\to$ \texttt{brcm,bcm2711-gpio}
        \item size of all GPIO registers? $\to$ \texttt{0xb4} (from \texttt{reg})
        \item find the gpiochip number for sysfs? $\to$
              \texttt{ls gpiochip*/device/driver}
        \item switch on the LED at GPIO19 (= sysfs nr \important{340}):
    \end{itemize}
\begin{lstlisting}[language=bash, style=basesmol, aboveskip=2pt, belowskip=2pt]
echo 340 > /sys/class/gpio/export
echo out > /sys/class/gpio/gpio340/direction
echo 1   > /sys/class/gpio/gpio340/value
\end{lstlisting}
    \begin{itemize}
        \item mmap the GPIO\_LVL register -- base address:
              legacy \texttt{0x7e200000} / ARM-view \texttt{0xfe200000}
        \item value for the code's \texttt{\#define}:
              \texttt{GPIO\_BASE\_ADDR = 0xfe200000}
    \end{itemize}
\end{example2}

\subsection{Exam Exercise (Probepr\"ufung 2021)}

\begin{example2}{Exam 2021 A3: GPIO access}
    A device-tree GPIO node is given; access the GPIOs.
    \tcblower
    \begin{itemize}
        \item \important{a)} driver name: \texttt{brcm,bcm2711-gpio}
        \item \important{b)} register block size: \texttt{0xb4} bytes
        \item \important{c)} find the gpiochip nr: \texttt{ls gpiochip*/device/driver}
        \item \important{d)} LED at GPIO19 = sysfs nr \important{340}:\\
              \texttt{echo 340 > /sys/class/gpio/export}\\
              \texttt{echo out > /sys/class/gpio/gpio340/direction}\\
              \texttt{echo 1 > /sys/class/gpio/gpio340/value}
        \item \important{e)} GPIO base addr: Legacy \texttt{0x7e200000},
              ARM-View \texttt{0xfe200000}
        \item \important{f)} \texttt{\#define GPIO\_BASE\_ADDR 0xfe200000}
        \item \important{g)} \texttt{gpio\_reg} (level reg, offset \texttt{0x34}) with
              \texttt{virtual\_gpio\_base = 0x84ec0000} $\to$
              \important{\texttt{0x84ec0034}}
    \end{itemize}
\end{example2}

\raggedcolumns
